% Generated draft: requires mathematical and source review.
\documentclass[11pt]{article}
\usepackage[margin=25mm]{geometry}
\usepackage{fontspec}
\setmainfont{Latin Modern Roman}
\usepackage{amsmath,amssymb,mathtools,unicode-math}
\setmathfont{Latin Modern Math}
\usepackage{enumitem,needspace,xurl,hyperref}
\setlength{\parindent}{0pt}
\setlength{\parskip}{5pt}
\setlength{\emergencystretch}{3em}
\newcommand{\N}{\mathbb{N}}
\newcommand{\Z}{\mathbb{Z}}
\newcommand{\Q}{\mathbb{Q}}
\newcommand{\R}{\mathbb{R}}
\newcommand{\C}{\mathbb{C}}
\newcommand{\F}{\mathbb{F}}
\newcommand{\A}{\mathbb{A}}
\newcommand{\PP}{\mathbb{P}}
\newcommand{\E}{\mathbb{E}}
\newcommand{\eps}{\varepsilon}
\newcommand{\dd}{\,\mathrm{d}}
\begin{document}
\section*{UnsolvedMath ID: 800584}
\textbf{Source problem code:} catalogue entry
\subsection{Definitions and mathematical statement}
The supplied record does not specify the intended class of spaces, define the \(\delta\)-property, or state whether covers must be countable. The following is a conditional formalization using the usual open-cover and limit-inferior interpretation; its agreement with the cited source requires review.

Let \(X\) be a topological space. An open cover is a family \(\mathcal U\) of proper open subsets of \(X\) whose union is \(X\). An \(\omega\)-cover is an open cover such that every finite subset \(F\subseteq X\) is contained in some \(U\in\mathcal U\). A \(\gamma\)-cover is an infinite open cover \(\mathcal V\) such that
\[
\forall x\in X,\qquad \{V\in\mathcal V:x\notin V\}\text{ is finite}.
\]
Write \(\Omega(X)\) and \(\Gamma(X)\) for these respective classes. If the source restricts admissible covers, these classes must be restricted consistently; no such restriction is specified in the supplied record.

The property \(\binom{\Omega}{\Gamma}\) means
\[
\forall\mathcal U\in\Omega(X)\ \exists\mathcal V\subseteq\mathcal U,
\qquad \mathcal V\in\Gamma(X).
\]
For a sequence \((A_n)_{n\geq1}\) of subsets of \(X\), define
\[
\liminf_{n\to\infty}A_n
=\bigcup_{m\geq1}\bigcap_{n\geq m}A_n.
\]
Let \(L(\mathcal U)\) be the smallest subfamily of \(\mathcal P(X)\) containing \(\mathcal U\) and closed under this operation on sequences of its members. In the interpretation adopted here, the \(\delta\)-property is
\[
\forall\mathcal U\in\Omega(X),\qquad X\in L(\mathcal U).
\]
The sets in \(L(\mathcal U)\) need not be open.

The source's question is whether, for every \(X\) in its intended class of spaces,
\[
X\text{ has the }\delta\text{-property}
\quad\Longleftrightarrow\quad
X\text{ satisfies }\binom{\Omega}{\Gamma}.
\]
The source states the implication from right to left, so the requested implication is from left to right. The intended class of spaces and cover conventions remain unresolved.

For context, \(\mathsf S_1(\Omega,\Gamma)\) means that for every sequence \((\mathcal U_n)_{n\geq1}\) of members of \(\Omega(X)\), there exist choices \(U_n\in\mathcal U_n\) such that \(\{U_n:n\geq1\}\in\Gamma(X)\). The supplied statement attributes the equality of properties \(\mathsf S_1(\Omega,\Gamma)=\binom{\Omega}{\Gamma}\) to reference [GN]. This is retained as supplied context, subject to the source's unresolved hypotheses and conventions.
\subsection{Short English statement}
Does every space with the \(\delta\)-property have a \(\gamma\)-subcover of each open \(\omega\)-cover? The source states the reverse implication. The formulation above uses a conventional interpretation of the undefined \(\delta\)-property; the intended spaces and cover restrictions require source review. The registry labels the problem open, but that status has not been independently verified.
\subsection{Sources}
\begin{itemize}
\item Local UnsolvedMath registry, supplied complete source record: permanent ID \(800584\), exact source problem code \texttt{catalogue entry}, title ``Selection Principles in Mathematics --- Problem 10 ([GN]).'' The record lists Tsaban as proposer, gives the year 2006, and labels the problem open.
\item Supplied source list: Tsaban, \emph{Selection Principles and Special Sets of Reals: Open Problems} (2006), Problem 10. Supplied URL: \url{https://arxiv.org/abs/math/0609601}. The record reports extraction by \texttt{arxiv\_latex} and an access date of 2026-10-16. This access information is supplied provenance, not an independently verified visit.
\item Reference label [GN] appears in both the attribution/context field and the asserted equality \(\mathsf S_1(\Omega,\Gamma)=\binom{\Omega}{\Gamma}\). The supplied record does not provide the corresponding bibliographic entry.
\end{itemize}
\end{document}
